\section{The regularized field and connecting paths}
\label{app:paths}

The finite-interval action argument requires a symmetric matrix path
from $K$ to $-K$ whose derivative measures the difference between a
matrix field and $AA^{\mathsf T}$.  We construct both the field and the
path here.  These are the constructions of \cite[Sections~4--5]{MLT};
we include their proofs to make the analytic input to the equality
argument self-contained.  The positivity statement below also
identifies the resolvent to which Theorem~\ref{trace:deficit} applies.

Throughout this appendix,
$K=\operatorname{diag}(k_1,\ldots,k_n)$ with $k_i>0$,
$\kappa=\max_i k_i$, and $\Sym_n$ denotes the real symmetric matrices.
We use $\beta=\sigma-\tfrac12$ and the constant $c_\sigma$ from
\eqref{trace:normalization}.

\subsection{A field with a scalar boundary integral}

For $\delta>0$, define on $y\geq0$
\begin{equation}\label{field:f}
 f_\delta(y)=\frac{c_\sigma}{\pi}\int_0^\infty t^\beta
       \left(\frac1{t+\delta}-\frac1{t+\delta+y}\right)\,dt,
\end{equation}
and extend to $y<0$ by the tangent line $f_\delta'(0)y$.
The resulting function is $C^1$, strictly increasing, and concave.
Indeed, for $y\geq0$,
\begin{equation}\label{field:f-derivative}
 f_\delta'(y)=a_\sigma(y+\delta)^{\beta-1},
 \qquad
 a_\sigma=\frac{c_\sigma}{\pi}
       \int_0^\infty\frac{u^\beta}{(1+u)^2}\,du>0.
\end{equation}
Both integrals converge since $-\tfrac12<\beta<\tfrac12$.
The extension allows the functional calculus below even when its
matrix argument has negative eigenvalues.

Set
\begin{equation}\label{field:definition}
 \begin{split}
 Q_s(B)&=s^2I-2sB+K^2,\\
 M_\delta(B)&=\lim_{R\to\infty}\int_{-R}^{R}
       \bigl(f_\delta(Q_s(B))-f_\delta(s^2)I\bigr)\,ds,
       \qquad B\in\Sym_n,
 \end{split}
\end{equation}
where the symmetric cutoff cancels the leading term
$-2s f_\delta'(s^2)B$, which is odd in $s$.
The scalar counterpart is
\begin{equation}\label{field:mu}
 \mu_\delta(z)=\int_\R
     \bigl(f_\delta(s^2+z)-f_\delta(s^2)\bigr)\,ds,
       \qquad z\in\R.
\end{equation}

\begin{proposition}[Regularized field]\label{field:properties}
The limit in \eqref{field:definition} exists locally uniformly, and
$M_\delta:\Sym_n\to\Sym_n$ is locally Lipschitz.  There is a
$C^1$ function $J_\delta:\Sym_n\to\R$, normalized by
$J_\delta(0)=0$, such that
\begin{equation}\label{field:primitive}
 dJ_\delta(B)[H]=-\tr(M_\delta(B)H).
\end{equation}
The scalar function $\mu_\delta$ is locally Lipschitz and strictly
increasing, with $\mu_\delta(0)=0$, and
\begin{equation}\label{field:mu-positive}
       \mu_\delta(z)=(z+\delta)^\sigma-\delta^\sigma
       \qquad(z\geq0).
\end{equation}
For every real unit vector $e$, and for every standard basis vector
$e_i$, respectively,
\begin{align}
 e^{\mathsf T}M_\delta(B)e
  &\leq\mu_\delta\bigl(e^{\mathsf T}K^2e
                         -(e^{\mathsf T}Be)^2\bigr),
       \label{field:direction}\\
 Be_i=ye_i\quad&\Longrightarrow\quad
 M_\delta(B)e_i=\mu_\delta(k_i^2-y^2)e_i.
       \label{field:coordinate}
\end{align}
For each diagonal sign matrix $S$,
\begin{equation}\label{field:equivariance}
                 M_\delta(SBS)=SM_\delta(B)S.
\end{equation}
Finally,
\begin{equation}\label{field:endpoints}
 \begin{split}
 J_\delta(-K)-J_\delta(K)
 &=\sum_{i=1}^n\int_{-k_i}^{k_i}\mu_\delta(k_i^2-s^2)\,ds\\
 &\longrightarrow b_\sigma\sum_{i=1}^n k_i^{2\sigma+1}
       \qquad(\delta\downarrow0).
 \end{split}
\end{equation}
\end{proposition}

\begin{proof}
We first justify the matrix integral, including its local Lipschitz
dependence.  If a scalar function $h$ is Lipschitz on an interval
containing the spectra of Hermitian matrices $X,Y$, then
\begin{equation}\label{field:HS-calculus}
 \|h(X)-h(Y)\|_{\HS}\leq\operatorname{Lip}(h)\|X-Y\|_{\HS}.
\end{equation}
To see this, use eigenbases $e_i$ of $X$ and $v_j$ of $Y$.
The mixed matrix entries of the two differences are
$(h(\lambda_i)-h(\nu_j))\langle e_i,v_j\rangle$ and
$(\lambda_i-\nu_j)\langle e_i,v_j\rangle$; square the scalar
Lipschitz inequality and sum over $i,j$.

On a bounded set of $B$'s the spectrum of
$E_s(B)=-2sB+K^2$ is contained in $[-a|s|,a|s|]$ for all sufficiently
large $|s|$.  For
\[
 h_s(y)=f_\delta(s^2+y)-f_\delta(s^2)-f_\delta'(s^2)y,
\]
Equation~\eqref{field:f-derivative} gives
$\sup_{|y|\leq a|s|}|h_s'(y)|=O(|s|^{2\beta-3})$.
It follows from \eqref{field:HS-calculus} that
\[
 f_\delta(Q_s(B))-f_\delta(s^2)I
      =-2s f_\delta'(s^2)B+T_s(B),
\]
where both $\|T_s(B)\|_{\HS}$ and the local Lipschitz constant of
$T_s$ are $O(|s|^{2\beta-2})$.  For the first bound, include
$f_\delta'(s^2)K^2$ in $T_s$; for the second, apply
\eqref{field:HS-calculus} to $h_s(E_s(B))$ and use
$E_s(B_1)-E_s(B_2)=-2s(B_1-B_2)$.
Because $2\beta-2<-1$, these remainders are integrable at infinity.
The odd term cancels at each symmetric cutoff.  On bounded
$s$-intervals, \eqref{field:HS-calculus} applies directly to
$f_\delta$.  This proves the asserted convergence and regularity.

For the primitive, choose a scalar antiderivative
$g_\delta'=f_\delta$.  Trace differentiation gives
$d\tr(g_\delta(X))[H]=\tr(f_\delta(X)H)$.
This formula follows for polynomials by cyclicity and for
$g_\delta$ by approximation in $C^1$ on compact spectral intervals.
For $s\ne0$, the function
\[
 B\longmapsto
 \frac{\tr(g_\delta(Q_s(B))-g_\delta(K^2))}{2s}
                  +f_\delta(s^2)\tr B
\]
has differential
$-\tr((f_\delta(Q_s(B))-f_\delta(s^2)I)H)$.
At $s=0$ that differential is constant in $B$ and also has a
primitive.  Thus the negative of every finite-cutoff field has
zero integral around each closed piecewise smooth curve in
$\Sym_n$.  Local uniform convergence passes this property to
$-M_\delta$.  Path integration, starting at zero, defines
$J_\delta$ and proves \eqref{field:primitive}.

For bounded $z$, the integrand in \eqref{field:mu} and its
Lipschitz constant in $z$ are $O(|s|^{2\beta-2})$ at infinity.
This proves absolute convergence and local Lipschitz continuity
of $\mu_\delta$.  Strict increase and $\mu_\delta(0)=0$ follow
from the same properties of $f_\delta$.  For $z\geq0$,
Tonelli's theorem applied to \eqref{field:f} gives
\[
 \mu_\delta(z)=c_\sigma\int_0^\infty t^\beta
       \bigl((t+\delta)^{-1/2}-(t+\delta+z)^{-1/2}\bigr)\,dt.
\]
The normalization \eqref{trace:normalization}, after scaling,
evaluates this as \eqref{field:mu-positive}.

Translations in the variable $s$ remain legitimate with a
symmetric cutoff.  More precisely, for fixed $a,z\in\R$,
\begin{equation}\label{field:shift}
 \lim_{R\to\infty}\int_{-R}^{R}
   \bigl(f_\delta((s-a)^2+z)-f_\delta(s^2)\bigr)\,ds
          =\mu_\delta(z).
\end{equation}
Indeed $F_z(s)=f_\delta(s^2+z)$ is even and
$F_z'(s)=O(|s|^{2\beta-1})\to0$ at infinity.  Pairing the two
short intervals produced by a translation bounds the difference
between the shifted and unshifted cutoff integrals by
\[
 |a|^2\sup_{R-|a|\leq u\leq R+|a|}|F_z'(u)|,
\]
which tends to zero.

Scalar concavity in an eigenbasis of $Q_s(B)$ now gives
\[
 e^{\mathsf T}f_\delta(Q_s(B))e
     \leq f_\delta(e^{\mathsf T}Q_s(B)e).
\]
The scalar argument on the right is
$(s-e^{\mathsf T}Be)^2+e^{\mathsf T}K^2e-(e^{\mathsf T}Be)^2$.
Integration and \eqref{field:shift} prove \eqref{field:direction}.
If $Be_i=ye_i$, functional calculus gives equality in the
corresponding coordinate, proving \eqref{field:coordinate}.
Since $SK^2S=K^2$, conjugation inside the defining integral proves
\eqref{field:equivariance}.

On diagonal matrices, \eqref{field:primitive} and
\eqref{field:coordinate} read
\[
 dJ_\delta(\operatorname{diag}(b_1,\ldots,b_n))
       =-\sum_i\mu_\delta(k_i^2-b_i^2)\,db_i.
\]
Integrating from $K$ to $-K$ gives the first equality in
\eqref{field:endpoints}.  For $z\geq0$,
$0\leq\mu_\delta(z)\leq z^\sigma$ and
$\mu_\delta(z)\to z^\sigma$.  Dominated convergence and the
substitution $s=k_i u$ give its stated limit.
\end{proof}

The field is defined on all symmetric matrices, but the action
argument compares it with a positive matrix $AA^{\mathsf T}$.
At such values its directional bound forces positivity of the
entire quadratic symbol $Q_s(B)$.  This is the condition needed
for the trace deficit formula.

\begin{proposition}[Positive values of the field]\label{field:positive}
If $M=M_\delta(B)\geq0$, then
\begin{equation}\label{field:positive-bounds}
 Q_s(B)\geq0\quad(s\in\R),\qquad
             0\leq M\leq\kappa^{2\sigma}I.
\end{equation}
Set $V=K^2-B^2+\delta I$ and use $B,V$ in
\eqref{trace:resolvents} and \eqref{trace:N}.  Then the symbol
$(sI-B)^2+V=Q_s(B)+\delta I$ is bounded below by $\delta I$, and
\begin{equation}\label{field:N}
                         N=M+\delta^\sigma I.
\end{equation}
In particular, the hypotheses of Theorem~\ref{trace:deficit} hold.
\end{proposition}

\begin{proof}
By \eqref{field:direction}, positivity of $M$ and strict increase
of $\mu_\delta$ imply
$e^{\mathsf T}K^2e\geq(e^{\mathsf T}Be)^2$ for every real unit
vector $e$.  Completing the scalar square gives
\[
 e^{\mathsf T}Q_s(B)e=(s-e^{\mathsf T}Be)^2
                +e^{\mathsf T}K^2e-(e^{\mathsf T}Be)^2\geq0.
\]
Thus $Q_s(B)\geq0$, also as a Hermitian matrix on $\C^n$.
Moreover, \eqref{field:direction} gives
$M\leq\mu_\delta(\kappa^2)I\leq\kappa^{2\sigma}I$, using
$(a+b)^\sigma\leq a^\sigma+b^\sigma$.

Write $Y_s=(sI-B)^2$.  Diagonalization of $B$ and
\eqref{field:shift} allow replacement of the scalar subtraction
$f_\delta(s^2)I$ in \eqref{field:definition} by $f_\delta(Y_s)$.
Since both $Q_s(B)$ and $Y_s$ are positive, their difference under
$f_\delta$ is, by \eqref{field:f},
\[
 \frac{c_\sigma}{\pi}\int_0^\infty t^\beta
 \left((Y_s+(t+\delta)I)^{-1}
          -(Q_s(B)+(t+\delta)I)^{-1}\right)\,dt.
\]
For fixed $B,\delta$, both matrices being inverted are bounded
below by $c(1+t+s^2)I$.  Their difference is $K^2-B^2$, so the
resolvent identity bounds the norm of the integrand by
$C t^\beta(1+t+s^2)^{-2}$.  Integrating first in $s$ yields
$C't^\beta(1+t)^{-3/2}$, which is integrable on $(0,\infty)$.
Fubini's theorem is therefore applicable and gives
\[
 M=c_\sigma\int_0^\infty t^\beta
        \bigl((t+\delta)^{-1/2}I-R(t)\bigr)\,dt.
\]
Subtracting this expression from \eqref{trace:N} and scaling
\eqref{trace:normalization} gives \eqref{field:N}.
\end{proof}

\subsection{Choosing the lower triangular coefficient matrix}

Let $\mathcal L_n$ be the real lower triangular $n\times n$
matrices.  Their dimension equals that of $\Sym_n$.  This dimension
match allows us to prescribe a symmetric terminal value by choosing
$C\in\mathcal L_n$.  At a simple initial problem there is one choice
up to signs of the rows.  Continuation preserves the parity of this
number, and hence produces the required terminal value.

\begin{proposition}[Connecting path]\label{path:connecting}
Let $\Omega=(x_-,x_+)$ be bounded and let
$U\in H^1_0(\Omega;\R^{n\times d})$, where $n,d\geq1$, satisfy
\[
                         \int_\Omega UU^{\mathsf T}\,dx=I_n.
\]
For every $\delta,\varepsilon>0$ there are
$C_\varepsilon\in\mathcal L_n$ and
$B_\varepsilon\in C^1(\overline\Omega;\Sym_n)$ such that, with
$A_\varepsilon=C_\varepsilon U$,
\begin{equation}\label{path:ode}
 \varepsilon B_\varepsilon'
       =M_\delta(B_\varepsilon)-A_\varepsilon A_\varepsilon^{\mathsf T},
 \qquad B_\varepsilon(x_-)=K,\quad B_\varepsilon(x_+)=-K.
\end{equation}
They satisfy
\begin{equation}\label{path:bounds}
 \sup_{x\in\overline\Omega}\|B_\varepsilon(x)\|_{\op}\leq\kappa,
 \qquad
 \sup_{0<\varepsilon\leq1}\|C_\varepsilon\|_{\HS}<\infty,
\end{equation}
where the second bound is for fixed $\delta,K,\Omega$.
\end{proposition}

\begin{proof}
Use the continuous representative of $U$.  Choose a smooth,
nonnegative, compactly supported cutoff $\chi$ on $\Sym_n$, equal
to one on $\{B:\|B\|_{\op}\leq\kappa\}$ and invariant under
conjugation by diagonal sign matrices.  Averaging a cutoff over
the finite sign group produces such a function.  The truncated
field $\widehat M=\chi M_\delta$ is bounded and globally Lipschitz.

Fix $\varepsilon>0$.  For $C\in\mathcal L_n$ and
$\theta\in[0,1]$, the equation
\begin{equation}\label{path:homotopy}
 \varepsilon B'=\theta\widehat M(B)-(CU)(CU)^{\mathsf T},
                         \qquad B(x_-)=K
\end{equation}
has a unique $C^1$ solution on $\overline\Omega$.  Define the
continuous endpoint map
\[
                  \Phi(C,\theta)=B(x_+)+K\in\Sym_n.
\]
Let $\mathcal G$ be the group of diagonal sign matrices, acting
by $C\mapsto SC$ on $\mathcal L_n$ and by conjugation on
$\Sym_n$.  Equivariance of the field and uniqueness of the ODE give
\begin{equation}\label{path:equivariance}
                       \Phi(SC,\theta)=S\Phi(C,\theta)S.
\end{equation}

We record the three facts needed for continuation: compactness
of the zero set, freeness of the sign action there, and one regular
orbit of zeros at $\theta=0$.
At a zero, integration of \eqref{path:homotopy} gives
\begin{equation}\label{path:C-identity}
       CC^{\mathsf T}=2\varepsilon K
                       +\theta\int_\Omega\widehat M(B(x))\,dx.
\end{equation}
Consequently all zeros lie in a bounded set, and their joint set
in $\mathcal L_n\times[0,1]$ is compact.  For
$0<\varepsilon\leq1$, this calculation gives the uniform bound
\begin{equation}\label{path:C-bound}
 \|C\|_{\HS}^2\leq2\tr K
       +n|\Omega|\sup_{B\in\Sym_n}\|\widehat M(B)\|_{\op}.
\end{equation}

No row of $C$ can vanish at a zero.  Otherwise the reflection
changing only row $i$ would fix $C$; ODE uniqueness and
\eqref{path:equivariance} would force $Be_i=y e_i$ throughout the
interval.  Equation~\eqref{field:coordinate} would then give
\[
 \varepsilon y'=\theta\chi(B(x))\mu_\delta(k_i^2-y^2),
                         \qquad y(x_-)=k_i.
\]
The coefficient $\chi(B(x))$ is continuous and the right side
is locally Lipschitz in $y$.  Uniqueness forces $y\equiv k_i$,
contradicting $y(x_+)=-k_i$.  Thus the action of $\mathcal G$
on the zero set is free.

At $\theta=0$ one has
\begin{equation}\label{path:initial-map}
                     \Phi(C,0)=2K-\varepsilon^{-1}CC^{\mathsf T}.
\end{equation}
Its lower triangular zeros are exactly
$C=\operatorname{diag}(d_1,\ldots,d_n)$ with
$d_i=\pm\sqrt{2\varepsilon k_i}$.  Indeed, the first row of
$CC^{\mathsf T}=2\varepsilon K$ fixes $|d_1|$ and forces all
other entries of the first column to vanish, and induction does
the same for subsequent columns.  These $2^n$ zeros form one
sign orbit.  At each, the derivative in a lower triangular
direction $H$ has diagonal entries
$-2\varepsilon^{-1}d_i H_{ii}$ and lower off-diagonal entries
$-\varepsilon^{-1}d_j H_{ij}$ for $i>j$.  It is therefore an
isomorphism $\mathcal L_n\to\Sym_n$.

We next justify the continuation using only smooth approximation
and the regular-value argument for parity
\cite[Sections~2--4]{Milnor1965}.  Suppose there were no zero at
$\theta=1$.  Compactness and freeness give a bounded invariant
open set $O\subset\mathcal L_n$ containing the projection of
every zero, whose closure lies where all rows are nonzero.
There is a positive lower bound for $\|\Phi\|_{\HS}$ on
\begin{equation}\label{path:gap-set}
       (\partial O\times[0,1])\ \cup\ (\overline O\times\{1\}).
\end{equation}
Reparametrize $\theta$ by a smooth map equal to zero near zero
and equal to one at one.  The resulting map $F$ agrees with
\eqref{path:initial-map} on an initial collar and retains the gap
on \eqref{path:gap-set}.

Extend $F$ constantly beyond the two ends of the parameter
interval, smooth by convolution, and splice to
\eqref{path:initial-map} in a smaller initial collar.
Averaging the approximation $\Psi_0$ in the form
\[
       \Psi(C,\theta)=\frac1{|\mathcal G|}
                  \sum_{S\in\mathcal G}S\Psi_0(SC,\theta)S
\]
gives a smooth equivariant approximation $\Psi$ that still equals
the initial map on that collar.  It can be made uniformly close
enough on $\overline O\times[0,1]$ to preserve the positive gap.

Freeness permits an equivariant perturbation with arbitrary target
directions.  At each $C_0\in\overline O$, choose a ball whose
distinct sign translates are disjoint, and a smooth bump $b$
supported in this ball with $b(C_0)=1$.  For $H\in\Sym_n$ the map
\[
                        V_H(C)=\sum_{S\in\mathcal G}b(SC)SHS
\]
is equivariant and satisfies $V_H(C_0)=H$.
Choose a basis of target matrices and then a finite collection of
such neighborhoods covering $\overline O$.  This gives finitely
many smooth equivariant maps $V_1,\ldots,V_N$ whose values span
$\Sym_n$ at every point of $\overline O$.

Let $h\geq0$ be smooth, vanish near zero, and be positive outside
the collar where $\Psi$ is the initial map.  Choose it so that
every zero of $h$ lies within that collar, and set
\[
       \Psi_\alpha(C,\theta)
           =\Psi(C,\theta)+h(\theta)\sum_{j=1}^{N}\alpha_jV_j(C).
\]
For $\alpha$ in a sufficiently small open ball, the gap on
\eqref{path:gap-set} persists.  At a zero of this family with
$0<\theta<1$, its full differential in $(C,\theta,\alpha)$
is onto: when $h(\theta)>0$ the parameter directions span the
target, and when $h(\theta)=0$ the derivative in $C$ is the
isomorphism computed for \eqref{path:initial-map}.
Its interior zero set is consequently a smooth manifold.
Sard's theorem \cite{Sard1942}, applied to projection of this
manifold onto the parameter ball, gives arbitrarily small
$\alpha$ for which zero is a regular value of
$(C,\theta)\mapsto\Psi_\alpha(C,\theta)$.
For completeness, if the full derivative is $(L_1,L_2)$ in
the $(C,\theta)$ and $\alpha$ directions, regularity of that
projection means that each $w$ admits $v$ with
$L_1v+L_2w=0$.  Hence $\operatorname{ran}L_2\subset
\operatorname{ran}L_1$, and surjectivity of $(L_1,L_2)$ implies
surjectivity of $L_1$.

For such an $\alpha$, the zero set in $O\times[0,1]$ is a
compact smooth one-dimensional manifold with boundary.  The gap
excludes boundary zeros except at $\theta=0$, where the initial
collar gives exactly the $2^n$ regular zeros above.  The free
sign action has a quotient that is again a compact
one-dimensional manifold with boundary: a neighborhood disjoint
from its nontrivial translates supplies the local interval or
half-interval chart.  All $2^n$ boundary points form one orbit,
so the quotient has exactly one boundary point.  This is
impossible, since a compact one-dimensional manifold is a finite
union of circles and closed intervals.  Therefore the original
map $\Phi(\cdot,1)$ has a zero.

It remains to remove the cutoff.  For the path at this zero and
any fixed real unit vector $e$, put $y=e^{\mathsf T}Be$.
Whenever $|y|>\kappa$, Equation~\eqref{field:direction} yields
\[
 e^{\mathsf T}M_\delta(B)e
       \leq\mu_\delta(e^{\mathsf T}K^2e-y^2)<0.
\]
Since $\chi\geq0$ and $(CU)(CU)^{\mathsf T}\geq0$, the ODE
implies $y'\leq0$ there.  Both endpoint values lie in
$[-\kappa,\kappa]$.  An excursion above $\kappa$ cannot start
from $\kappa$ while $y$ is nonincreasing; an excursion below
$-\kappa$ cannot return to $-\kappa$ while $y$ is nonincreasing.
The two endpoint conditions therefore rule out both excursions.
This holds for every $e$, so $\|B(x)\|_{\op}\leq\kappa$.
Consequently $\chi(B)=1$, giving \eqref{path:ode}, and
\eqref{path:C-bound} supplies the remaining assertion in
\eqref{path:bounds}.
\end{proof}
